\section{Introduction}
\label{sec:introduction}

For the squared-distance cost on a compact Riemannian manifold,
positive upper and lower density bounds give an almost-everywhere
optimal map, but continuity depends on the geometry of the cost.
The weak Ma--Trudinger--Wang condition controls that geometry where
the cost is smooth. We prove that it gives a common H\"older estimate
for the map and its inverse throughout each fixed density-bounded
class, even when transport graphs approach conjugate pairs.

Let $(M,g)$ be a smooth connected compact Riemannian manifold without
boundary, of dimension $n\geq2$. Write $d$ and $\vol$ for its distance
and volume measure, and set
\[
 c(x,y)=\frac12d(x,y)^2,\qquad
 I(x)=\{p\in T_xM:\ d(x,\exp_x(ap))=a|p|
                         \text{ for some }a>1\}.
\]
We assume the weak Ma--Trudinger--Wang condition: for every
$p\in I(x)$ and $\xi,\eta\in T_xM$ with $\xi\perp\eta$,
\begin{equation}\label{eq:weak-mtw}
 -\frac32\left.\frac{\partial^4}{\partial s^2\partial t^2}
 c\bigl(\exp_x(t\xi),\exp_x(p+s\eta)\bigr)\right|_{s=t=0}\geq0.
\end{equation}
The cost is smooth near the pair in this formula. No curvature
inequality at a nonsmooth cost pair is assumed.
For probability densities $\rho_0,\rho_1$, an optimal map is a
measurable map $T$ with $T_\#(\rho_0\vol)=\rho_1\vol$ that minimizes
$\int_M c(x,T(x))\rho_0(x)\,d\vol(x)$ among maps with that pushforward.

\begin{theorem}[Uniform bi-H\"older transport]\label{thm:main}
Fix $(M,g)$ as above and $0<\lambda\leq\Lambda<\infty$ for which the
following density class is nonempty. There are
$\alpha\in(0,1]$ and $C<\infty$, depending only on
$(M,g),\lambda,\Lambda$, such that the following holds for every pair
of measurable densities satisfying
\[
 \int_M\rho_i\,d\vol=1,\qquad
 \lambda\leq\rho_i\leq\Lambda\quad\vol\text{-almost everywhere},
 \qquad i=0,1.
\]
The almost-everywhere unique optimal map from $\rho_0\vol$ to
$\rho_1\vol$ for the cost $c$ has a homeomorphic representative
$\widetilde T:M\to M$, and
\[
 d(\widetilde T(x),\widetilde T(x'))\leq C d(x,x')^\alpha,
 \qquad
 d(\widetilde T^{-1}(y),\widetilde T^{-1}(y'))
       \leq C d(y,y')^\alpha
\]
for all $x,x',y,y'\in M$.
\end{theorem}

The constants are common to the full density class on the fixed
manifold. They do not depend on individual densities or their
derivatives. The theorem includes conjugate cut points, and does not
assume nonfocality, strict MTW, or cut avoidance. We obtain a positive
exponent by contradiction; no explicit exponent or uniformity over
varying metrics is asserted.

\subsection{Context and geometric input}

McCann's polar factorization theorem supplies the almost-everywhere
unique optimal map and its cost potential \cite{McCann2001}.
The curvature condition originates in the work of
Ma--Trudinger--Wang \cite{MTW2005}; Loeper established its connection
with supporting cost functions and continuity \cite{Loeper2009}.
Sections and contact-volume comparison are central tools in
Caffarelli's theory \cite{Caffarelli1992} and in the weak-MTW
regularity theory of Figalli--Kim--McCann
\cite[Theorem~2.1 and Corollary~9.6]{FKM2013Holder}. These results already treat rough densities.
Their smooth-cost H\"older estimates require the appropriate
cost-convexity and nondegeneracy hypotheses. A density-controlled
exponent after localization does not by itself provide a common
localization scale or H\"older constant for a family whose transport
graphs approach conjugate pairs.

Uniform control over a density-bounded class is known in important
geometric settings. Loeper--Villani obtain such estimates under
strict MTW and nonfocality, with a further extension under their
uniform-regularity hypothesis and positivity of the minimizing-fiber
diameter $\delta(M)$
\cite[Corollaries~8.3--8.4]{LoeperVillani2010}.
Figalli--Kim--McCann prove uniform cut separation on products of
round spheres \cite[Theorem~5.1 and Corollary~5.2]{FKM2013Spheres}.
The issue addressed here is the general fixed weak-MTW manifold,
without imposing such separation. A different quantitative stability
result of Warren \cite[Theorem~1]{Warren2026} assumes common
$C^2$ norm bounds on the logarithms of smooth densities and sufficiently small
Wasserstein distance; those assumptions are not part of the
measurable-density class considered here.

The density-free input is proved in the companion paper
\emph{Global Support and Convex Injectivity Domains under Weak MTW}
\cite{Foundation2026}. It derives convexity of open and closed
injectivity domains, global support for every ordinary subgradient,
and uniform intermediate regularity from \eqref{eq:weak-mtw}.
We state the complete interface in Section~\ref{sec:geometry}.
The global supporting conclusion has an important antecedent:
Figalli--Rifford--Villani \cite[Lemma~3.1]{FRV2011} prove the
global double-mountain principle assuming convex injectivity
domains as well as weak MTW. The companion supplies the geometric
hypothesis under weak MTW alone, together with the intermediate
statements needed below. Figalli--Gallou\"et--Rifford
\cite[Theorem~1.7]{FGR2015} previously proved the convexity
implication under nonfocality. On compact surfaces,
Figalli--Rifford--Villani \cite[Theorem~1.3]{FRV2011} proved that
transport continuity is equivalent to the conjunction of convex
injectivity domains and weak MTW. Their introduction poses the
equivalence with weak MTW alone; the separate convexity and continuity
questions also appear in \cite[Questions~8.6--8.7]{LPZ2019}.

Figalli--Rifford--Villani's transport continuity property concerns
all density pairs bounded above and away from zero. They prove that
this property implies convex injectivity domains and weak MTW
\cite[Definition~1.1 and Theorem~1.2(i)]{FRV2011}.
Theorem~\ref{thm:main} implies that property by choosing bounds for
each pair. Combined with their necessity result, it yields the
weak-MTW--transport-continuity equivalence, with the forward
direction strengthened to the class-uniform estimates above.

\subsection{Conventions}

We use dual potentials with
\[
 u(x)=\sup_y\{-c(x,y)-v(y)\},\qquad
 v(y)=\sup_x\{-c(x,y)-u(x)\}.
\]
Thus $u(x)+c(x,y)+v(y)\geq0$, with equality at a contact.
The symbol $\partial u(x)$ denotes the ordinary subdifferential
of the semiconvex function, with covectors identified with tangent
vectors by the metric. At a differentiable contact its gradient
is the initial minimizing log. We write
\[
 G_x=\{p\in T_xM:d(x,\exp_xp)=|p|\}=\overline{I(x)},\qquad
 c_t=c/t,\qquad Q_tf(z)=\min_x\{f(x)+c_t(x,z)\},
\]
and, for $p\in I(x)$,
\[
 A_x(p)=D^2_{xx}c(x,\exp_xp),\qquad
 \sigma_x(p)=|\det(d\exp_x)_p|.
\]
The function $\sigma_x$ extends continuously to $G_x$. At an
ordinary cut point, a specified nonconjugate minimizing branch
has a smooth upper cost branch; we state explicitly when such a
branch is used. A star denotes the adjoint. Constants denoted by
$c,C$ may change between formulas. Constants in the final
comparison depend only on the fixed geometry and density bounds;
constants used solely to pass limits at fixed scales may also
depend on the already fixed templates.

\subsection{Proof structure}

For a dual pair $(u,v)$, the lifted gap section at $x$ and
height $r\geq0$ is
\[
 \mathcal S_x(r)=\left\{p\in G_x:
       u(x)+\frac{|p|^2}{2}+v(\exp_xp)\leq r\right\}.
\]
We study the oscillation of $v$ on its endpoint image
$\exp_x\mathcal S_x(r)$. Convexity of the lifted section
turns a power bound on that oscillation into a power bound
on the diameter of the endpoint image. Taking the oscillation
supremum over all admissible density pairs makes this a single
scalar estimate for the entire class. Once it is proved, every
zero-height section has one endpoint, and comparison of nearby
poles gives the H\"older estimate in both transport directions.

The foundation theorem applies to all ordinary subgradients, including
convex combinations of all minimizing logs at active cut endpoints.
The projection $(x,p)\mapsto\exp_x(tp)$ of the full subgradient graph
is a homeomorphism for $0<t<1$; only the supporting inequality is
asserted at time one for an arbitrary potential. It also gives compact
convex lifted gap sections and $C^{1,1}$ intermediate potentials.
The minimizing point in $Q_tu(z)$, called its pole, depends
Lipschitzly on $z$.

Section~\ref{sec:spectral} develops the quantitative estimates needed
when a minimizing endpoint approaches conjugacy. A fixed joining
point factors the full Jacobi matrix into nonsingular piece
matrices and a positive semidefinite broken-action Hessian.
Radial shortening adds a uniformly positive form. Transverse MTW,
the radial identity, and a common separating logarithmic scale
then compare ordered singular values along an extended affine
segment.

Section~\ref{sc:section} takes a supremum of section oscillations
over the entire density class. Failure of every power bound gives
two separated gap scales on which this supremum is almost constant.
Two scalar modifications produce a positive comparison maximum.
One modifies the value at a center endpoint; the other supplies
the active endpoints whose minimizing logs have the center log
as a convex combination.
Varying their amplitude gives an inequality for every positive
active barycentric representation, with constants fixed before
the curvature of the outer modification is chosen.

Section~\ref{sec:jets} keeps the order of the limiting argument
explicit. An exact pole-matrix identity transfers lower jets to
active endpoints. The original contact-volume bound first excludes
conjugacy at a selected outer endpoint and bounds its determinant.
This bounds the center Hessian. Only then does the short-time
minimizer have nonsingular derivative, allowing the original
almost-everywhere density equation to be used at the sampled
center endpoints. An oblique-projection determinant gain and the
spectral comparison give the contradiction. Finally,
Section~\ref{sec:conclusion} converts the resulting power bound
into the common estimate for the map and its inverse.
